\section{A two-interface inertial formation diagnostic}
\label{sec:reduced-formation}
The following is a deliberately restricted forward model, not a replacement
for the viscous, streaming, thermal and three-dimensional problem above.
Assume axisymmetric, incompressible, irrotational mean motion, impermeable
slip walls, and small interface displacement relative to each layer depth.
Evaluate the kinetic energy on the flat reference domains. Retain the
nonlinear graph capillary energy and recompute harmonic radiation loads on
the instantaneous graphs. No viscous damping or settling law is introduced.

Here is a direct derivation of the inertial coupling. For a radial mode
$J_0(kr)$, with $J_1(kR)=0$, the potential in each layer satisfies
$\phi_{zz}-k^2\phi=0$. The zero mode is excluded by the sealed-volume
constraints. If the bottom and top vertical velocities of a layer of depth
$h$ are $u$ and $v$, respectively, its potential is
\[
 \phi(z)=\frac{v-u\cosh(kh)}{k\sinh(kh)}\cosh(kz)
             +\frac{u}{k}\sinh(kz),\qquad 0\le z\le h.
\]
Integration of $\rho|\nabla\phi|^2/2$, or Green's identity, gives kinetic
energy per radial modal norm
\[
 \frac{\rho}{2k}\bigl[\coth(kh)(u^2+v^2)-2uv\,\operatorname{csch}(kh)\bigr].
\]
Consequently, the two-interface modal mass is
\[
 M_k=\frac1k\begin{pmatrix}
 \rho_0 c_0+\rho_1 c_1&-\rho_1 s_1\\
 -\rho_1 s_1&\rho_1 c_1+\rho_2 c_2
 \end{pmatrix}.
\]
Here $c_j=\coth(kh_j)$ and $s_j=\operatorname{csch}(kh_j)$.
Projection of the pinned, volume-preserving spline velocities into these
radial modes yields a symmetric positive mass matrix $M$. This projection
must be converged independently of the acoustic mesh. An optional Ritz
restriction retains the lowest eigenvectors of $K_0 e=\omega^2 M e$;
it is a declared dynamical resolution, not a source optimization.

With spline coordinates $q$, or their Ritz restriction, the forward equation is
\[
 M\ddot q+\nabla E(q)=F_{\rm ac}(q,a(t)g),\qquad q(0)=\dot q(0)=0.
\]
The target does not occur on the right-hand side. A specified source envelope
$a(t)$ is a physical command, not interpolation towards a target. With
zero drive, the continuous reduced model conserves $\dot q^TM\dot q/2+E(q)$.
For nonzero drive its derivative is $\dot q^T F_{\rm ac}$; no convergence
to an equilibrium is implied. A midpoint mechanical update with a predicted
midpoint acoustic solve is second order locally and must be checked at
fixed commands by time-step refinement. Large deformation, omitted modes,
unresolved wave fields and resonant sensitivity can invalidate this diagnostic.
Neither its error curve nor its optical rendering establishes the ten-nanometre
criterion of the full apparatus.
